\section{Introduction}
\label{sec:intro}

LLM agents have emerged as a new class of database client. Unlike human analysts, who issue queries irregularly as analytical intuition develops, agents follow systematic task-decomposition patterns: they establish a working dataset, then issue a sequence of aggregations, filters, and breakdowns against it. Across a recent characterization of agent sessions on analytical workloads, we observe that the same plan-critical subexpressions recur at strikingly high rates within a session, and that the agent frequently telegraphs its next several queries in its chain of thought before executing them. This regularity is not incidental---it is a direct consequence of how modern LLM agents are trained and prompted to decompose analytical tasks.

This regularity is an optimization opportunity that existing cooperative database mechanisms are poorly equipped to exploit. Materialized views, result recyclers~\cite{ivanova2010recycler}, adaptive indexing~\cite{idreos2007cracking}, and classical multi-query optimization~\cite{sellis1988mqo,roy2000mqo} all share a common design assumption: that reuse opportunities must be \emph{inferred} from observed workload history, after enough statistics have accumulated to justify the materialization or rewriting decision. This assumption is fundamental---it governs what the database is allowed to do, when, and at what cost. Under it, a database is reactive by necessity: it can only discover reuse after paying for it at least once.

LLM agents invalidate this assumption. An agent operating under a cooperation protocol can \emph{declare} its intent prospectively: which results will be reused, which queries will follow, which physical access patterns will dominate. Given such declarations, a far richer space of optimizations becomes \emph{legal}---not merely heuristic. A result that the agent declares will be reused can be materialized on first execution without waiting for recurrence to prove the value. A set of queries the agent declares it will issue next can be co-optimized as a batch, their shared subexpressions materialized once, their physical properties synthesized from the union of their access patterns. None of this is safe under inferred intent; all of it is safe under declared intent.

\paragraph{Extending reuse across the session.}
Modern query engines already provide \emph{intra-query} reuse for free. PostgreSQL 16's multi-reference CTE auto-materialization, for instance, detects common subexpressions \emph{within} a single query and materializes them once. What has been missing is \emph{inter-query} reuse---sharing materialized state across the sequence of queries that constitutes an agent's session. Concord provides exactly this extension: a session-scoped workspace, addressable by declared names, whose contents persist across the agent's query sequence and are optimized using the agent's declared reuse and lookahead intent.

\paragraph{Three cooperating mechanisms.}
Concord consists of three mechanisms, each enabled by a specific declaration:

\begin{itemize}
\item \textbf{M1: Declared materialization.} The agent emits \texttt{-- WILL\_SAVE} before a query whose result will be reused, or \texttt{-- WILL\_SAVE\_CTE} to mark a named CTE body within a larger query as the reusable fragment. The session harness materializes the declared result on first execution via \texttt{CREATE TEMP TABLE AS}, eliminating the double-execution cost of a post-hoc save.

\item \textbf{M2: Lookahead-informed multi-query optimization.} The agent emits \texttt{-- LOOKAHEAD} declaring its next $k$ queries. The harness identifies common subexpressions across the declared window, co-materializes them once with physical properties (sort order, partial aggregates, covering columns) synthesized from the union of declared access patterns, and rewrites the declared queries to exploit the shared result.

\item \textbf{M3: Reuse-aware physical tuning.} The harness adapts session-level GUCs (e.g., \texttt{work\_mem}, parallelism) to the declared workspace's characteristics, and chooses clustering and partial-aggregate options from the declared access pattern.
\end{itemize}

Concord is implemented as a session harness over \emph{unmodified} PostgreSQL 16: no extensions, no planner patches. The declared-intent protocol is delivered through SQL comments and a small PL/pgSQL workspace catalog. This deliberately minimal integration surface is itself a design claim: if the protocol requires deep database modifications, its deployability story is weak.

\paragraph{Contributions.}
We make four contributions:

\begin{itemize}
\item A declared-intent protocol for cooperative agent-database optimization, with explicit legality conditions stating what each declaration entitles the database to do that it could not safely do otherwise (Section~\ref{sec:protocol}).

\item An empirical characterization of agent adoption and predictability on IMDB/JOB: \adoptionRate voluntary save rate across 20 multi-step analytical sessions with no explicit reward for saving, \mTwoPrecision M2 declaration precision, and a detailed picture of \emph{when} agents declare intent faithfully versus \emph{when} they iterate on the base definition in ways that defeat reuse (Section~\ref{sec:motivation}, \ref{sec:eval}).

\item A cost-model result establishing that lookahead-informed shared materialization is never worse than ad-hoc per-query optimization under stated assumptions on declaration fidelity (Section~\ref{sec:protocol}).

\item An end-to-end evaluation showing Concord yields \imdbSpeedupFull session speedup on IMDB/JOB and \ibSpeedup on InsightBench at equal analytical quality, with the speedup net of PostgreSQL 16's existing intra-query CTE materialization (Section~\ref{sec:eval}).
\end{itemize}

\paragraph{Paper organization.}
Section~\ref{sec:related} situates Concord against four prior cooperative-database traditions. Section~\ref{sec:motivation} presents the workload characterization that motivates the protocol, including a negative result on cardinality-feedback optimization for agent workloads that clarifies why the opportunity lies at the session level rather than at the single-query level. Section~\ref{sec:protocol} formalizes the protocol, the three mechanisms, and the cost model. Section~\ref{sec:impl} describes the implementation over PostgreSQL 16. Section~\ref{sec:eval} reports the empirical evaluation. Section~\ref{sec:discussion} discusses failure modes and generalization; Section~\ref{sec:conclusion} concludes.
